\documentclass[10pt,letterpaper]{article}
\usepackage[utf8]{inputenc}
\usepackage[spanish]{babel}
\usepackage{listings}
\usepackage{graphicx}



\title{Entregable fase 4: Informe de validaci\'on y optimizaci\'on de tienda online}
\author{Shomara Rivera}
\date{\today}

\usepackage{xcolor}

\definecolor{codegreen}{rgb}{0,0.6,0}
\definecolor{codegray}{rgb}{0.5,0.5,0.5}
\definecolor{codepurple}{rgb}{0.58,0,0.82}
\definecolor{backcolour}{rgb}{0.95,0.95,0.92}

\lstdefinestyle{mystyle}{
    backgroundcolor=\color{backcolour},
    commentstyle=\color{codegreen},
    keywordstyle=\color{magenta},
    numberstyle=\tiny\color{codegray},
    stringstyle=\color{codepurple},
    basicstyle=\ttfamily\footnotesize,
    breakatwhitespace=false,
    breaklines=true,
    captionpos=t,
    keepspaces=true,
    numbers=left,
    numbersep=5pt,
    showspaces=false,
    showstringspaces=false,
    showtabs=false,
    tabsize=2,
    extendedchars=true,
    literate={á}{{\'a}}1 {é}{{\'e}}1 {í}{{\'i}}1 {ó}{{\'o}}1 {ú}{{\'u}}1 {Á}{{\'A}}1 {É}{{\'E}}1 {Í}{{\'I}}1 {Ó}{{\'O}}1 {Ú}{{\'U}}1 {ñ}{{\~n}}1 {Ñ}{{\~N}}1 {¿}{{\textquestiondown}}1 {¡}{{\textexclamdown}}1
}
\lstset{style=mystyle}

 \renewcommand{\lstlistingname}{C\'odigo}
 \renewcommand{\lstlistlistingname}{Lista de c\'odigos}

\begin{document}
\maketitle
\tableofcontents
\lstlistoflistings

\newpage

\section{Resultados de consultas con datos de prueba}
Al realizar las tres consultas del entregable anterior usando los datos de prueba que añad\'i, pudimos ver que el sistema responde de forma correcta y nos devuelve las tablas que esperamos; esto nos sirve para asegurar que la estructura soporta las operaciones de la tienda.

\begin{enumerate}
    \item Listar productos disponibles por categor\'ia, ordenados por precio:
    Esta consulta combina los productos con sus respectivas categor\'ias usando un \texttt{JOIN}, mostrando de menor a mayor los art\'iculos que s\'i tienen inventario.

    \begin{lstlisting}[language=SQL, frame=single, caption={Consulta de productos disponibles}]
SELECT p.nombre AS producto, c.nombre AS categoria, p.precio, p.stock
FROM productos p
JOIN categorias c ON p.id_categoria = c.id_categoria
WHERE p.stock > 0
ORDER BY p.precio ASC;
    \end{lstlisting}

    \begin{verbatim}

+--------------------+------------+----------+-------+
| producto           | categoria  | precio   | stock |
+--------------------+------------+----------+-------+
| Mica de Cristal    | Accesorios |   150.00 |   120 |
| Cable USB-C        | Accesorios |   200.00 |   100 |
| Funda iPhone 13    | Accesorios |   300.00 |    60 |
| Cargador Rápido    | Accesorios |   400.00 |    80 |
| Mouse Inalámbrico  | Accesorios |   500.00 |    45 |
| PowerBank 10000mAh | Accesorios |   800.00 |    30 |
| Teclado Mecánico   | Accesorios |  1200.00 |    20 |
| Galaxy Buds        | Accesorios |  3000.00 |    40 |
| AirPods Pro        | Accesorios |  5000.00 |    50 |
| Poco X3            | Teléfonos  |  5000.00 |    40 |
| Sony WH-1000XM4    | Accesorios |  6000.00 |    15 |
| Moto G100          | Teléfonos  |  7500.00 |    30 |
| Xiaomi Mi 11       | Teléfonos  |  8000.00 |    25 |
| Huawei P40         | Teléfonos  |  9000.00 |     8 |
| Pixel 6            | Teléfonos  | 10000.00 |    10 |
| OnePlus 9          | Teléfonos  | 11000.00 |     5 |
| Galaxy S21         | Teléfonos  | 12000.00 |    15 |
| Galaxy S22         | Teléfonos  | 14000.00 |    12 |
| Acer Swift         | Laptops    | 15000.00 |    15 |
| iPhone 13          | Teléfonos  | 15000.00 |    20 |
| Asus ZenBook       | Laptops    | 17000.00 |     9 |
| HP Envy            | Laptops    | 18000.00 |    12 |
| iPhone 14          | Teléfonos  | 18000.00 |    10 |
| MacBook Air        | Laptops    | 20000.00 |    10 |
| ThinkPad X1        | Laptops    | 22000.00 |     7 |
| Lenovo Legion      | Laptops    | 24000.00 |     6 |
| Dell XPS 13        | Laptops    | 25000.00 |     5 |
| MSI Stealth        | Laptops    | 28000.00 |     4 |
| MacBook Pro        | Laptops    | 30000.00 |     3 |
| Razer Blade        | Laptops    | 35000.00 |     2 |
+--------------------+------------+----------+-------+
    \end{verbatim}

    \item Mostrar clientes con pedidos pendientes y total de compras:
    Aqu\'i agrupamos los pedidos con estado pendiente para identificar qu\'e usuarios tienen compras en proceso, el conteo se hace de forma individual gracias al \texttt{GROUP BY}.

    \begin{lstlisting}[language=SQL, frame=single, caption={Consulta de pedidos pendientes}]
SELECT c.nombre AS cliente, c.correo, COUNT(pe.id_pedido) AS total_pedidos_pendientes
FROM clientes c
JOIN pedidos pe ON c.id_cliente = pe.id_cliente
WHERE pe.estado = 'pendiente'
GROUP BY c.id_cliente;
    \end{lstlisting}

    \begin{verbatim}
+-----------------+---------------------+---------------------------+
| cliente         | correo              | total_pedidos_pendientes  |
+-----------------+---------------------+---------------------------+
| Angel Duenas    | angel@email.com     |                         2 |
| Paco Luna       | paco@email.com      |                         1 |
| Nolan Rivera    | nolan@email.com     |                         1 |
| José Ruz        | jose@email.com      |                         1 |
| Jhovanny Pérez  | jhovanny@email.com  |                         1 |
+-----------------+---------------------+---------------------------+
\end{verbatim}

    \item Reporte de los 5 productos con mejor calificaci\'on promedio en rese\~nas:
    Esta consulta calcula el promedio de estrellas usando la funci\'on \texttt{AVG}; el filtro \texttt{LIMIT 5} recorta la tabla para mostrarnos solo el top de los mejores.

    \begin{lstlisting}[language=SQL, frame=single, caption={Consulta de promedio de calificaciones}]
SELECT p.nombre AS producto, AVG(r.calificacion) AS promedio_calificacion
FROM productos p
JOIN resenas r ON p.id_producto = r.id_producto
GROUP BY p.id_producto
ORDER BY promedio_calificacion DESC
LIMIT 5;
    \end{lstlisting}

    \begin{verbatim}
+-------------------+-----------------------+
| producto          | promedio_calificacion |
+-------------------+-----------------------+
| iPhone 13         |                5.0000 |
| Mouse Inalámbrico |                5.0000 |
| Galaxy S21        |                5.0000 |
| Cargador Rápido   |                5.0000 |
| Dell XPS 13       |                5.0000 |
+-------------------+-----------------------+
    \end{verbatim}
\end{enumerate}


\section{Explicaci\'on de \'indices creados y su impacto}
Como vimos en las diapositivas de clase, para mejorar el rendimiento de la base de datos agregamos \'indices en las columnas que m\'as vamos a usar para buscar datos, como un \'indice funciona como el \'indice de un libro, permite que MySQL vaya directo a la fila en lugar de leer toda la tabla.

\begin{itemize}
    \item \textbf{idx\_producto\_nombre:} Creado sobre el nombre de los productos, ayuda a que el buscador de la tienda sea m\'as r\'apido.
    \item \textbf{idx\_producto\_categoria:} Creado sobre la llave for\'anea de categor\'ias para filtrar r\'apido.
    \item \textbf{idx\_pedido\_cliente:} Creado en los pedidos para cargar el historial del cliente sin demoras.
\end{itemize}

Para comprobar que s\'i funcionan, utilizamos el comando \texttt{EXPLAIN} antes de hacer un SELECT, pues este comando le pregunta a MySQL c\'omo planea buscar el dato.

\begin{verbatim}
+----+-------------+-----------+------+---------------------+---------------------+---------+-------+------+----------+
| id | select_type | table     | type | possible_keys       | key                 | key_len | ref   | rows | filtered |
+----+-------------+-----------+------+---------------------+---------------------+---------+-------+------+----------+
|  1 | SIMPLE      | productos | ref  | idx_producto_nombre | idx_producto_nombre | 602     | const |    1 |   100.00 |
+----+-------------+-----------+------+---------------------+---------------------+---------+-------+------+----------+
\end{verbatim}

Al ver la columna llamada \texttt{key}, podemos ver que MySQL s\'i utiliz\'o nuestro \'indice; y en la columna \texttt{rows}, dice que solo tuvo que leer 1 fila para encontrar el resultado, demostrando que la consulta est\'a optimizada.


\section{Pruebas de procedimientos almacenados (escenarios de validaci\'on y error)}
Para asegurar que se cumplan las restricciones requeridas, ejecutamos los procedimientos almacenados usando el comando \texttt{CALL}, intentando meter tanto datos falsos para ver si el sistema nos deten\'ia, como datos correctos para verificar las actualizaciones.

\subsection{Intento de rese\~na sin compra previa}
Intentamos que el cliente 2 hiciera una rese\~na del producto 1, el cual nunca ha comprado.
\begin{lstlisting}[language=SQL, frame=single, caption={Error por falta de compra}]
CALL sp_registrar_resena(5, 'Excelente celular', 1, 2);
\end{lstlisting}
El procedimiento detect\'o que no hab\'ia registros de compra y bloque\'o la inserci\'on, arrojando el mensaje: \textit{Error: Solo puedes hacer rese\~nas de productos que hayas comprado.}

\subsection{Intento de compra sin stock suficiente}
Intentamos registrar un pedido por 1000 unidades de un producto que solo tiene 20 en stock.
\begin{lstlisting}[language=SQL, frame=single, caption={Error por falta de stock}]
CALL sp_registrar_pedido(1, 1, 1000, 15000.00);
\end{lstlisting}
El procedimiento us\'o un IF para revisar el stock y, al ver que no alcanzaba, abort\'o la venta arrojando: \textit{Error: No hay suficiente stock para este producto.}

\subsection{Intento de agregar producto duplicado}
Intentamos registrar un "iPhone 13" cuando ya existe otro igual en la misma categor\'ia.
\begin{lstlisting}[language=SQL, frame=single, caption={Error por producto duplicado}]
CALL sp_agregar_producto('iPhone 13', 'Tel\'efono de prueba', 15000.00, 10, 1);
\end{lstlisting}
El sistema cont\'o los nombres repetidos y lanz\'o la alerta: \textit{Error: Ya existe un producto con ese nombre en esta categor\'ia.}

\subsection{Actualizaci\'on del inventario de productos}
Probamos el procedimiento encargado de restarle art\'iculos al stock simulando que se vendieron 5 piezas del producto con ID 2.
\begin{lstlisting}[language=SQL, frame=single, caption={Actualizaci\'on exitosa de stock}]
CALL sp_actualizar_stock(2, 5);
\end{lstlisting}
Al ejecutar el \texttt{UPDATE}, el sistema redujo el stock de forma correcta sin necesidad de modificar los dem\'as campos del producto.

\subsection{Cambio en el estado log\'istico de un pedido}
Simulamos que el pedido n\'umero 1 cambi\'o de estatus en el proceso de entrega de la tienda.
\begin{lstlisting}[language=SQL, frame=single, caption={Modificaci\'on del estado del pedido}]
CALL sp_cambiar_estado_pedido(1, 'enviado');
\end{lstlisting}
MySQL modific\'o la columna estado de 'pendiente' a 'enviado' gui\'andose estrictamente por el ID proporcionado.

\subsection{Eliminaci\'on de una rese\~na y c\'alculo del promedio}
Mandamos a borrar la rese\~na n\'umero 1 para verificar que el sistema vuelva a calcular las estrellas del producto.
\begin{lstlisting}[language=SQL, frame=single, caption={Eliminaci\'on de rese\~na y recalculo de promedio}]
CALL sp_eliminar_resena(1);
\end{lstlisting}
El sistema aplic\'o el \texttt{DELETE} y de inmediato nos mostr\'o en pantalla la tabla con el nuevo promedio de calificaciones corregido.

\subsection{Actualizaci\'on de los datos de contacto de un cliente}
Modificamos el perfil del cliente n\'umero 1 simulando que cambi\'o su n\'umero telef\'onico y su casa.
\begin{lstlisting}[language=SQL, frame=single, caption={Actualizaci\'on de perfil de cliente}]
CALL sp_actualizar_cliente(1, 'Avenida Principal 123', '555-4321');
\end{lstlisting}
El sistema reemplaz\'o la direcci\'on y el tel\'efono viejos, dejando intactos el nombre y el correo electr\'onico del usuario.

\subsection{Generaci\'on del reporte automatizado de stock bajo}
Invocamos la consulta de control para extraer la lista de productos que tienen menos de 5 piezas en el almac\'en.
\begin{lstlisting}[language=SQL, frame=single, caption={Llamada al reporte de stock bajo}]
CALL sp_reporte_stock_bajo();
\end{lstlisting}
Al no requerir variables de entrada, se llama con los par\'entesis vac\'ios y MySQL Workbench nos arroja la lista ordenada de los art\'iculos que urge resurtir.


\section{Propuestas de mejoras adicionales}
Al ir haciendo este proyecto me di cuenta de un par de cosas que podr\'iamos implementar a futuro bas\'andonos en los temas que vimos en clase:

\begin{itemize}
    \item \textbf{Uso de vistas (VIEWS):} Durante la Fase 3 me di cuenta de que hacer consultas con \texttt{JOIN} puede ser un poco complejo, como mejora, propongo guardar las consultas m\'as dif\'iciles (como el reporte de clientes con pedidos pendientes) dentro de una VIEW, as\'i, cualquier usuario podr\'ia ver la informaci\'on r\'apidamente sin tener que escribir todo el c\'odigo SQL cada vez.
    \item \textbf{Uso de triggers (Disparadores):} Para tener a\'un m\'as seguridad, propongo agregar un Trigger antes de insertar rese\~nas, aunque el procedimiento ya revisa que el cliente haya comprado el art\'iculo, un Trigger nos ayudar\'ia a asegurar que nadie meta un n\'umero negativo o un 10 en la calificaci\'on de estrellas, garantizando que siempre sea del 1 al 5.
\end{itemize}


\end{document}